\section{Considerações Finais}

\subsection{Síntese do Atendimento aos Requisitos do Projeto}
A Tabela \ref{tab:requisitos} sintetiza o atendimento formal aos requisitos técnicos e práticos estabelecidos para esta etapa do projeto.

\begin{table}[htbp]
\centering
\caption{Síntese do Atendimento aos Requisitos do Projeto}
\label{tab:requisitos}
\scriptsize
\setlength{\tabcolsep}{3pt}
\begin{tabular}{lp{4.2cm}c}
\toprule
\textbf{Item} & \textbf{Requisito Técnico} & \textbf{Status} \\
\midrule
\textbf{Item 1} & Configuração do ambiente (Quartus e ferramentas). & 100\% \\
\textbf{Item 2} & Caminho de dados com saltos (\texttt{j}, \texttt{jal}, \texttt{jr}). & 100\% \\
\textbf{Item 3} & Validação de todos os módulos via testbenches. & 100\% \\
\textbf{Item 4} & Parametrização integral (\texttt{parameter}). & 100\% \\
\textbf{Item 5} & Apresentação em vídeo e relatório no padrão IEEE. & 100\% \\
\bottomrule
\end{tabular}
\end{table}

\subsection{Conclusão e Trabalhos Futuros}
O desenvolvimento da Segunda Avaliação consolidou a implementação de um processador MIPS Monociclo de 32 bits em Verilog HDL. O hardware incorporou suporte às instruções \texttt{sll}, \texttt{slti} e \texttt{bne} para viabilizar a execução do algoritmo \texttt{sort.bin}. Além disso, implementou-se o mascaramento de barramento na memória ROM e a inicialização de registradores no reset para contornar o modelo de memória virtual do simulador MARS.

A validação automatizada comprovou a ordenação do vetor na memória de dados com zero falhas através de 7 bancadas de teste. A síntese lógica no Intel Quartus Prime para o FPGA Intel MAX 10 concluiu sem latches espúrios e com total coerência estrutural.

Na Terceira Avaliação Parcial o caminho de dados será segmentado nos cinco estágios clássicos de pipeline: busca, decodificação, execução, memória e escrita. Essa etapa introduzirá registradores de desacoplamento temporal, unidade de detecção de conflitos e unidade de adiantamento de operandos.
